In special cases, we use the notation $U_F = 0$ and $SU_F = SO$ for $F = \mathbf{R}$, $U_F = U$ and $SU_F = SU$ for $F = \mathbf{C}$, and $U_F = Sp$ for $F = \mathbf{H}$.

We recall that $V_k(F^n)$ is the subspace of $F^{kn}$ consisting of $k$-tuples $(v_1, \ldots, v_k)$, where $(v_i|v_j) = \delta_{i,j}$ for $1 \leq i, j \leq k$. Here we may have $n = \infty$, in which case $V_k(F^\infty)$ has the inductive topology defined by the following subsets:

$$V_k(F^k) \subset V_k(F^{k+1}) \subset \cdots \subset V_k(F^n) \subset \cdots \subset V_k(F^\infty) = \bigcup_{k \leq n} V_k(F^n)$$

We define a map $\eta_k^n$ or just $\eta$: $U_F(n) \rightarrow V_k(F^n)$ by the relation $\eta(u) = (u(e_1), \ldots, u(e_k))$. The map $\eta$ is a continuous surjection. Moreover, the following diagram is commutative.

$$\begin{array}{ccc}
